\chapter{Constructing a causal model and identifying its questions}
\label{r3:graphs}
\section{Begin with the intervention, unit and observation process}
Consider a household eligible for a preferential deposit product. Let $Z$
denote the bank's invitation, $A$ actual activation, $M$ a subsequent change
in balances and $Y$ the ninety-day funding contribution. An intervention on
$Z$ is operationally different from forcing $A=1$. An increase in the
database field recording $A$ is different again. The bank can randomize an
invitation; it cannot generally compel activation or create liquidity merely
by changing a recorded balance. A causal question must identify the mechanism
that the proposed decision actually changes.

The unit is a household rather than an account when money can move between
accounts owned by the same people. The outcome includes an agreed transfer
price, interest paid and servicing cost, with movements within the household
consolidated. This is part of defining $Y$, rather than a correction performed
after estimating an account-level effect. For a fixed population and horizon,
write $Y(a)$ for the outcome under a specified activation intervention, and
$Y^{Z=z}$ for the outcome under invitation assignment $z$. Two initial targets
are
\begin{equation}
 \theta_A=\E\{Y(1)-Y(0)\},\qquad
 \theta_Z=\E\{Y^{Z=1}-Y^{Z=0}\}.
 \label{r3:targets}
\end{equation}
They differ even in an ideal randomized invitation experiment. An invitation
may have costs or affect trust without activation. Such effects belong to
$\theta_Z$ and threaten an exclusion restriction used to infer $\theta_A$.

Consistency means that a household exposed to the specified action has its
corresponding potential outcome. It requires clarity about price, duration,
channel, eligibility, execution and outcome measurement. It is not ensured
by giving all those interventions the same binary code. Interference would
require $Y_i(\mathbf a)$, or a justified exposure mapping, instead of $Y_i(a_i)$.
The present single-period derivations assume a population of noninterfering
units and a fixed version of each action. Cluster or network designs must
change those assumptions explicitly \citep[Chapters 1--3]{hernan2020}.

For clarity most graph calculations below use discrete variables and sums.
Continuous versions use conditional densities or probability kernels, with
absolute continuity and integrability stated at the estimation stage. This
distinction matters: an intervention at one exact continuous treatment value
can be identifiable as a conditional-law functional without admitting the
same regular root-$n$ estimator as a binary-treatment average effect.

\section{Structural equations define what an intervention changes}
\begin{definition}[Acyclic structural model]
An acyclic structural causal model $\mathcal M$ consists of endogenous
variables $V=(V_1,\ldots,V_d)$, exogenous variables $U$, a joint law $P_U$,
and assignments
\[
 V_j:=f_j(V_{\pa(j)},U_j),\qquad j=1,\ldots,d,
\]
whose endogenous parent graph is a DAG. A hard intervention $\doop(A=a)$
replaces the assignment for $A$ by the constant $a$ and retains the remaining
assignments and exogenous law.
\end{definition}
The colon distinguishes an assignment from an algebraic identity that may
be inverted freely. A regression residual is constructed to be orthogonal
to fitted regressors; a structural disturbance represents omitted causes
and need not be orthogonal to the action. Autonomous mechanisms are what
make intervention replacement meaningful \citep[Sections 3.1--3.2]{pearl2009}.

If all $U_j$ are mutually independent, the model is Markovian relative to
its displayed endogenous DAG. Hidden common causes can instead be displayed
as latent parents, producing a larger DAG, or left as dependence in $P_U$.
One must not draw an observed DAG and silently assume independent disturbances
after omitting a common cause. Both the functions and the disturbance law
are components of the causal model.

Potential outcomes can be generated by recursively evaluating the modified
assignments at the same exogenous realization: $Y(a)=f_Y(\ldots;a,U)$.
This constructs consistency and expresses $P\{Y(a)\}$ as the distribution
of $Y$ in the intervened model. Conversely, a potential-outcome analysis
need not posit mutually independent disturbances for every structural
equation. In particular, assumptions sufficient to identify each
single-world intervention distribution need not determine joint quantities
involving incompatible interventions.

The SWIG construction makes this distinction explicit. Split an action node
into its random part $A$, retaining incoming arrows, and a fixed part $a$,
retaining outgoing arrows. Relabel its descendants as counterfactual
variables. For $X\to A$, $X\to Y$, $A\to Y$ this yields a graph for
$(X,A,Y(a))$ in which $A\perp Y(a)\mid X$ can be read by separation.
It does not simultaneously supply a graph for the unobserved pair
$(Y(0),Y(1))$. The SWIG framework covers the less restrictive FFRCISTG
model as well as the independent-error structural submodel
\citep[Section 2, node splitting and modularity]{swig2013}. Consequently,
identifying a mean effect does not identify individual gains or their
variance: $\Var\{Y(1)-Y(0)\}$ additionally involves the unobserved
$\Cov\{Y(1),Y(0)\}$.

\section{A diagram is constructed from mechanisms}
\label{r3:graph-construction}
Before drawing an arrow, define the decision clock. Here $X$ is measured
before invitation, $Z$ is assigned next, $A$ records subsequent activation,
$M$ is later behavior, and $Y$ closes the horizon. Let $U$ be unrecorded
liquidity need and $V$ a servicing difficulty that also affects outcomes.
Let $C$ indicate whether the episode has a complete outcome record. A
candidate mechanism graph is shown in Figure~\ref{r3:bank-graph}. It is
an assumption to be challenged with process owners and data provenance.

\begin{figure}[htbp]\centering
\begin{tikzpicture}[>=Stealth,node distance=15mm,
 every node/.style={draw,circle,minimum size=8mm},latent/.style={dashed,draw,circle}]
 \node (Z) at (0,0) {$Z$};
 \node (A) at (2,0) {$A$};
 \node (M) at (4,0) {$M$};
 \node (Y) at (6,0) {$Y$};
 \node (X) at (3,2) {$X$};
 \node[latent] (U) at (3,-1.6) {$U$};
 \node (C) at (2,-3.1) {$C$};
 \node[latent] (V) at (6,-3.1) {$V$};
 \draw[->] (Z)--(A); \draw[->] (A)--(M); \draw[->] (M)--(Y);
 \draw[->] (A) to[bend left=15] (Y);
 \draw[->] (X)--(A); \draw[->] (X)--(M); \draw[->] (X)--(Y);
 \draw[->] (U)--(A); \draw[->] (U)--(Y);
 \draw[->] (A)--(C); \draw[->] (V)--(C); \draw[->] (V)--(Y);
\end{tikzpicture}
\caption{A proposed deposit-activation model. Dashed nodes are unobserved.
The absence of $Z\to Y$ encodes an exclusion restriction, while selection
on a complete record $C=1$ can open $A\to C\leftarrow V\to Y$.}
\label{r3:bank-graph}
\end{figure}

Each omitted arrow makes a substantive claim. Absence of $Z\to Y$ and
$Z\to M$ says that the invitation affects the economic outcome only through
activation. Absence of causes shared by $Z$ and $(U,X,V)$ represents genuine
random assignment in this example. Absence of $U\to M$ says that liquidity
need affects later balances only through activation or through mechanisms
already represented by $X$; this might be implausible and should be revised.
A discovered small regression coefficient cannot establish any of these
omissions. Arrows allow a mechanism; their presence does not prescribe its
sign, magnitude, linearity or monotonicity.

A useful construction procedure has five concrete steps.
\begin{enumerate}
 \item Define the actor, feasible action, unit, economic endpoint and clock.
       Separate assigned, received, recorded and measured quantities.
 \item List causes of the action and of the endpoint, including private
       decisions, calendar conditions, common counterparties and selection.
 \item Order mechanisms at the finest relevant time resolution. Replace
       apparent contemporaneous feedback by a justified within-period order
       or an explicitly cyclic/dynamic model.
 \item Mark what is observed and when it becomes available; draw observation
       and inclusion mechanisms instead of treating missingness as a blank cell.
 \item Examine every exclusion needed by the intended identification route.
       Record its business justification, plausible violations and what an
       experiment or additional measurement could resolve.
\end{enumerate}
This procedure separates a causal diagram from data lineage. A lineage edge
from transactions to an aggregated balance describes a calculation. It does
not mean that intervening on the stored aggregate changes the underlying
transactions. Deterministic accounting identities also create conditional
dependencies that require care in discovery algorithms relying on generic
faithfulness assumptions.

\section{Paths, colliders, latent variables and learned graphs}
\label{r3:separation}
A path ignores the direction of its successive edges. An internal node is
a collider when both adjacent edges have arrowheads into it. A DAG path is
active conditional on $S$ when every noncollider is outside $S$ and every
collider has itself or a descendant in $S$. Two sets are d-separated when
no such path connects them. In a distribution Markov to the DAG,
d-separation implies conditional independence
\citep[Definition 1]{pearl2009}. Faithfulness is the additional converse
assumption used by many discovery procedures; it is not needed to deduce
independences from a stipulated causal graph.

In the bank graph, conditioning on $X$ blocks $A\leftarrow X\to Y$ but
leaves $A\leftarrow U\to Y$ open. Adding $M$ can block a causal path and
therefore change the target. Conditioning on $A$ opens the collider
$Z\to A\leftarrow U\to Y$: invitation and latent need need not remain
independent among activators. This is why analyzing compliance strata defined
by observed receipt is not the same as randomizing receipt. Conditioning on
$C=1$ can also open $A\to C\leftarrow V\to Y$.

For a numerical collider example, let $A,V$ be independent fair binary
variables, $Y=V$, and let $C=1$ exactly when $A=V$. Action has zero causal
effect on $Y$. Among records with $C=1$, however, $Y=A$, so the observed
contrast is one. Positive marginal action probabilities and perfect outcome
prediction within the selected sample do not repair selection bias.

Marginalizing latent nodes can yield an acyclic directed mixed graph (ADMG).
In the simple common-cause motif, $A\leftarrow U\to Y$ becomes
$A\leftrightarrow Y$; a directed latent path can induce a directed observed
edge. Directed and bidirected edges can coexist. The corresponding separation
rule, m-separation, uses arrowheads to define colliders in the same way.
A bidirected edge records a permitted hidden common-cause route; it is not
an instruction to regress in two directions. With selection and more general
latent structures, maximal ancestral graphs and PAGs require their own
semantics; an ADMG is not interchangeable with those objects. The
latent-time-series setting of \citet[Section 2]{gerhardus2020}, for example,
explicitly excludes selection
and targets a time-window-dependent ancestral equivalence object.

A DAG supplied by an analyst, a CPDAG representing a DAG Markov-equivalence
class, and a PAG representing latent/selection-compatible ancestral graphs
are different inputs. For DAGs the same skeleton and unshielded colliders
give the same implied conditional independences
\citep[Definition 6.24 and Lemma 6.25]{peters2017}. Two variables alone do
not observationally distinguish $A\to Y$ from $Y\to A$. Known timing may
exclude the latter but still leave hidden confounding. A discovery score
cannot turn the remaining graph ambiguity into a unique causal effect by
choosing whichever orientation permits an estimator to run.

\section{Identification is a property of a model class}
\label{r3:identification}
Let $\mathfrak M$ be the class of causal models permitted by the assumptions,
$P_{\mathcal M}$ their observed-data law, and $\theta(\mathcal M)$ the
requested effect. Identification means
\begin{equation}
 P_{\mathcal M_1}=P_{\mathcal M_2}
 \quad\Longrightarrow\quad
 \theta(\mathcal M_1)=\theta(\mathcal M_2)
 \quad\text{for every }\mathcal M_1,\mathcal M_2\in\mathfrak M.
 \label{r3:id-definition}
\end{equation}
It permits writing $\theta=\Theta(P)$ for a statistical functional
\citep[Definition 2]{pearl2009}. Identification does not assert that the
functional is well conditioned, estimable at root-$n$ rate, or supported
with useful precision by the observed sample. Conversely, failure to find
an adjustment set does not prove nonidentification: another transformation
may identify the target.

\begin{proposition}[Truncated factorization for a fully observed Markovian DAG]
Assume independent structural disturbances, acyclicity and autonomous
assignments. For a hard intervention on a node set $I$, the density of the
remaining variables is
\begin{equation}
 p^a(v_{-I})=\prod_{j\notin I}p(v_j\mid v_{\pa(j)})\big|_{v_I=a}.
 \label{r3:truncated}
\end{equation}
Observed conditional kernels identify the right-hand side wherever the
required parent configurations have support.
\end{proposition}
\begin{proof}
Order nodes topologically. Independent disturbances imply that each
assignment generates a kernel conditional on its parents, giving the
observational product $\prod_j p(v_j\mid v_{\pa(j)})$. Replacing an
assignment by a constant removes its kernel and substitutes its value in
the arguments of each child kernel. The other kernels remain invariant
by autonomy. Iterating in the same order gives Equation~\eqref{r3:truncated}.
Without support, a conditional kernel is not determined by observational
data at the parent configurations that the intervention creates.
\end{proof}
This is the causal Markov and truncated-factorization argument in
\citet[Theorem 1 and Corollary 1]{pearl2009}. For a graph with hidden $U$,
the full product may still be written and integrated over $U$, but its
factors involving $U$ are not automatically observable. Replacing them by
conditionals of observed variables alone is an additional identification
claim, not a consequence of the chain rule.

\subsection{Adjustment and its potential-outcome proof}
Suppose a pretreatment set $X$ blocks all paths entering $A$ through an
incoming arrow and contains no descendants of $A$. The back-door criterion
is a sufficient graphical route to exchangeability in the stated model;
it is not the most general adjustment criterion
\citep[Definition 3]{pearl2009}. In potential-outcome notation the needed
conditions for each action $a$ are consistency, $Y(a)\perp A\mid X$, and
$P(A=a\mid X=x)>0$ for almost every target $x$. Then
\begin{align}
 \E\{Y(a)\}
 &=\sum_x\E\{Y(a)\mid X=x\}p(x)\nonumber\\
 &=\sum_x\E\{Y(a)\mid A=a,X=x\}p(x)\nonumber\\
 &=\sum_x\E[Y\mid A=a,X=x]p(x).
 \label{r3:adjustment}
\end{align}
The first step averages over the target population; the second uses
exchangeability; the third uses consistency, with positivity making the
conditional mean observable. The graph is one justification for those
equalities, not an alternative to them. For the bank graph containing
$A\leftarrow U\to Y$, $X$ alone does not justify the second step.

If instead all common causes are measured in $X$ and the mediator has the
law $p(m\mid a,x)$, the total-effect mean is
\[
 \sum_{x,m}\E[Y\mid a,m,x]p(m\mid a,x)p(x).
\]
The mediator distribution must change with $a$. Regressing on $M$ and
reading off an action coefficient while holding $M$ fixed generally
answers a different question. The difference is visible before choosing
a regression family.

\subsection{Front-door identification when adjustment fails}
Consider the distinct graph $A\to M\to Y$ with a hidden common cause of
$A$ and $Y$, no direct $A\to Y$, no confounding of $A\to M$, and all
back-door paths from $M$ to $Y$ blocked by $A$. Assume the required
$p(a,m)$ support. Under these restrictions,
\begin{equation}
 p(y\mid\doop(a))
 =\sum_m p(m\mid a)\sum_{a'}p(y\mid m,a')p(a').
 \label{r3:frontdoor}
\end{equation}
To see the steps, $A\to M$ is unconfounded, so its interventional law is
$p(m\mid a)$. Setting $M=m$ cuts the incoming $A\to M$ edge; adjustment
by $A$ identifies $p(y\mid\doop(m))$ as the inner sum. In this graph
intervening on $A$ changes $Y$ only by changing $M$; the response kernel
under an intervention on $M$ is invariant to that $A$ intervention.
Composing these two intervention kernels gives the displayed expression.
It is not ordinary adjustment for a post-treatment covariate
\citep[Section 6.7, Example 6.46]{peters2017}.

For deposits, measured balance change is rarely a defensible front-door
mediator: liquidity need may directly affect it and funding contribution
may have activation-related cost channels outside balances. The formula
is useful because it states which measurements and exclusions would be
needed. It must not motivate deleting inconvenient arrows from
Figure~\ref{r3:bank-graph}.

Do-calculus generalizes these transformations by conditions for deleting
observations, exchanging actions and observations, and deleting actions
in modified graphs. For example,
\[
 p(y\mid\doop(a),z,x)=p(y\mid\doop(a),x)
 \quad\text{if }Y\perp_G Z\mid A,X\text{ in }G_{\bar A},
\]
where $G_{\bar A}$ removes edges entering $A$ and the separation statement
is evaluated in that modified graph. The action/observation exchange rule
requires a different modification, removing outgoing edges of the variable
being exchanged. A compiler must retain the actual rule and separation
witness for each step. The next subsection explains how an identification
algorithm organizes such operations and distinguishes a failed search
from a proof of nonidentification.

\subsection{District decomposition and an obstruction certificate}
\label{r3:id-algorithm}
In a semi-Markovian ADMG a district is a maximal set joined by bidirected
paths. A singleton without a bidirected neighbor is also a district.
For $S\subseteq V$, define the interventional kernel
$Q[S]=p(s\mid\doop(v\setminus s))$. It is a function of $s$ and the
fixed outside values; in general it is not the ordinary conditional
$p(s\mid v\setminus s)$. Shared hidden disturbances couple variables
within a district. Integrating over distinct independent latent groups
factorizes an intervention kernel over districts in its remaining graph.

The ID algorithm exploits this decomposition recursively. With no
intervention it marginalizes the supplied law. It first restricts to
ancestors of the requested outcome, summing irrelevant variables. It
can then add interventions on remaining nonancestors in the graph with
the original intervention nodes removed, because they cannot change the
target in that graph. When the remaining graph has several districts,
it identifies their kernels separately, multiplies them, and marginalizes
nonoutcome variables. When one district remains, it either extracts
topologically ordered conditional factors for a district already
available in the current kernel, recurses inside a larger containing
district, or returns a failure witness. These steps are made precise in
\citet[Section 4, Figure 2]{id2006}; they apply to its specified
semi-Markovian model class.

The obstruction is a \emph{hedge}. An $R$-rooted C-forest is a
bidirected-connected subgraph whose observed nodes have at most one
directed child, with childless root set $R$. Two such forests
$F'\subset F$ form a hedge for the query when $F$ meets the intervention
set, $F'$ does not, and their common roots are ancestors of the outcome
in the intervened graph. The source proves that such a pair witnesses
nonidentification and establishes completeness of ID for the specified
queries. These are inherited theorems; the following small case makes
their practical meaning concrete.

For the bow graph $A\to Y$, $A\leftrightarrow Y$, take
$F=\{A,Y\}$ and $F'=\{Y\}$, with common root $Y$. The larger
forest contains $A$ and the smaller does not. Thus failure to identify
$p(y\mid\doop(a))$ is certified, rather than merely caused by a
search routine finding no adjustment set. Section~\ref{r3:nonidentification}
constructs two positive observed-law examples witnessing that same
failure. The front-door graph avoids this obstruction for the total
effect through an identifiable composition of district kernels.

Conditional interventional queries require another distinction. Dividing
an identified $p(y,z\mid\doop(a))$ by its $z$ margin works when that
joint law is identified and the denominator is positive, but failure
to identify the joint law need not settle the conditional query.
The IDC algorithm first exchanges eligible conditioning variables for
interventions using do-calculus, then invokes ID and normalizes the
remaining joint kernel \citep[Section 5, Figure 3]{id2006}.
Its input is a conditional interventional distribution, not a generic
cross-world mediation or personalized counterfactual request.

\subsection{A recursive call trace, including a conditional rescue}
\label{r3:id-trace}
The front-door functional can now be reconstructed as an execution of
ID, rather than supplied as a memorized template. Use the topological
order $(A,M,Y)$ in the three-node front-door graph above. The original
districts are $\{A,Y\}$ and $\{M\}$. Write $\mathrm{ID}(B,I;K,G)$
for a call identifying the law of $B$ after setting $I$, from the supplied
kernel $K$ and graph $G$. The line numbers below refer to
\citet[Figure 2]{id2006}.
\begin{enumerate}
 \item In $\mathrm{ID}(Y,A;p,G)$, all nodes are ancestors of $Y$.
 Removing $A$ leaves districts $\{M\}$ and $\{Y\}$. Line 4 returns
 $\sum_m\mathrm{ID}(M,\{A,Y\};p,G)
             \mathrm{ID}(Y,\{A,M\};p,G)$.
 \item For the first factor, ancestor reduction (line 2) restricts the
 graph to $\{A,M\}$ and the input to $p(a,m)$. The intervention set
 becomes $\{A\}$. In this graph $\{M\}$ is an original district;
 line 6 extracts $p(m\mid a)$.
 \item For the second factor, the remaining singleton $\{Y\}$ is
 contained in the larger original district $\{A,Y\}$. Line 7 recurses
 inside that district with intervention $A$ and kernel
 \[
    q(a,y\mid m)=p(a)p(y\mid a,m).
 \]
 This is a normalized kernel with $m$ fixed. It is generally different
 from $p(a,y\mid m)$ because $p(a\mid m)\ne p(a)$.
 \item The induced graph on $\{A,Y\}$ has a bidirected edge and no
 directed edge. Its ancestors of $Y$ are just $\{Y\}$. Line 2
 therefore marginalizes $a$, leaving
 $q_Y(y\mid m)=\sum_{a'}p(a')p(y\mid a',m)$ and an empty
 intervention set. Line 1 returns that kernel.
 \item Multiplication and the final sum over $m$ recover
 Equation~\eqref{r3:frontdoor}. The dummy $a'$ in the second kernel
 is distinct from the intervention value $a$ in the first.
\end{enumerate}
For a fully specified finite check, let $U$ be fair binary and use
independent innovations with probabilities
\[
 P(A=1\mid U)=.2+.6U,\quad P(M=1\mid A)=.1+.7A,\quad
 P(Y=1\mid M,U)=.1+.5M+.2U.
\]
All observed cells are positive. Direct intervention gives
$\E[Y(a)]=.25+.35a$, while the observed regression gives
$\E[Y\mid A=a]=.19+.47a$. Enumerating the latent law and evaluating
the ID expression agree on the effect $.35$; direct adjustment for no
covariates produces $.47$. The numerical check validates this supplied
example and its algebra, not an implementation of general ID.

A second graph shows why IDC has a separate role. Let
$A\to M\to Y$ and $A\leftrightarrow M$, with no other edges.
The joint $p(m,y\mid\doop(a))$ is not identifiable: marginalizing
an identified joint would identify $p(m\mid\doop(a))$, contradicted
by the bow-graph hedge on $\{A,M\}$. Yet the conditional is identified.
In $G_{\bar A,\underline M}$, where incoming edges to $A$ and outgoing
edges from $M$ are removed, $Y$ is separated from $M$ given $A$.
IDC line 1 exchanges the observation of $M$ for an action:
\[
 p(y\mid\doop(a),m)=p(y\mid\doop(a),\doop(m)).
\]
ID then extracts the original singleton district $\{Y\}$, giving
$p(y\mid a,m)=p(y\mid m)$. Positivity is still required at the
conditioning value. Normalizing a failed joint-identification result
would have missed this valid conditional query. Conversely, the result
does not identify the distribution of $M$ under the action or the
unconditional action effect. This trace exemplifies the distinction in
\citet[Figure 3 and Theorems 7--8]{id2006}.

An identification compiler should therefore carry typed kernels and
intervention sets through each transformation. Multiplying a district
kernel as if it were an ordinary regression conditional is a semantic
error even if the expression is numerically evaluable. A graph outside
the supported class should produce an unsupported-method result;
an actual hedge or paired-model witness supports nonidentification.

\section{Two observationally identical worlds}
\label{r3:nonidentification}
Here is a strictly positive counterexample adapted as a project derivation
to the identification definition. In model I, $U$ is a fair binary variable,
$P(A=1\mid U)=0.2+0.6U$, and, independently of the action noise,
$P(Y=1\mid U,A)=0.2+0.6U$. Action has no effect. In model II, $A$ is
independent and fair, with $P(Y=1\mid A)=0.32+0.36A$. Both yield
\[
 \begin{array}{c|cc}
  &Y=0&Y=1\\\hline
 A=0&0.34&0.16\\
 A=1&0.16&0.34
 \end{array}
\]
because, for instance, model I gives
$P(A=1,Y=1)=\tfrac12(0.2^2+0.8^2)=0.34$. Their ATEs are zero and
$0.36$, respectively. A class allowing hidden $A/Y$ confounding and an
action effect contains both models. Its observed law therefore fails
Equation~\eqref{r3:id-definition}. More data can estimate the common
table accurately but cannot choose between the intervention distributions.

One can still describe an identified set. For binary $Y$, under consistency
alone, let $q_a=P(Y=1,A=a)$ and $p_a=P(A=a)$. The unobserved part of
$\E[Y(a)]$ lies between zero and $1-p_a$, so
\begin{equation}
 \theta_A\in[q_1-q_0-p_1,\ q_1+p_0-q_0].
 \label{r3:consistency-bounds}
\end{equation}
These elementary bounds follow by assigning the missing binary potential
outcomes to their extremes; simultaneous assignments attain the endpoints
under the unrestricted consistency model. For the table above they are
$[-0.32,0.68]$. They do not encode an IV, monotonicity or substantive
confounding sensitivity restriction. Adding one of those assumptions means
optimizing over a smaller, explicitly stated model class.

\section{Selection, mediation and the population being answered}
Suppose $C=1$ marks observed outcomes. If $Y\perp C\mid A,X$ and
$\pi(a,x)=P(C=1\mid A=a,X=x)>0$, and $X$ and $A$ are observed for the
full target cohort, then
\[
 \E[Y\mid a,x]=\E\!\left[\frac{CY}{\pi(a,x)}\,\middle|\,a,x\right].
\]
The identity follows by conditioning on $Y,A,X$ and canceling $\pi$.
It can be combined with Equation~\eqref{r3:adjustment} only if the causal
adjustment assumptions also hold. If servicing difficulty $V$ in
Figure~\ref{r3:bank-graph} affects both completeness and outcome, this
conditional observation assumption may fail. Conditioning on retention
or survival does not automatically define an effect for the original
eligible population. When an endpoint is undefined after an event, an
imputation rule alone cannot make the causal target well defined.

Mediation adds another distinction. A controlled direct effect compares
$\E[Y(1,m)-Y(0,m)]$ at a feasible fixed mediator intervention. A natural
direct effect compares $\E[Y(1,M(0))-Y(0,M(0))]$. The latter combines
counterfactuals from incompatible worlds and requires restrictions beyond
those for the total effect. Under consistency, suitable support, and
assumptions implying $Y(a,m)\perp M(a')\mid X$ together with the
required action and mediator exchangeabilities, one obtains
\[
 \E[Y(a,M(a'))]
 =\sum_{x,m}\E[Y\mid a,m,x]p(m\mid a',x)p(x).
\]
The cross-world independence is an additional assumption, not something
proved by fitting the two observed regressions. Treatment-induced
mediator--outcome confounding invalidates this simple construction. A
single-world analysis and the stronger independent-error model must
therefore remain distinguishable \citep{swig2013}.

\subsection{Edge assignments and a single-world alternative}
\label{r3:mediation-interventions}
For the simple independent-error system
$M=f_M(A,X,U_M)$ and $Y=f_Y(A,M,X,U_Y)$, with conditional
independence of the relevant disturbances given $X$, an edge
intervention can send $a'$ along $A\to M$ and $a$ along $A\to Y$.
The resulting response is
$f_Y(a,f_M(a',X,U_M),X,U_Y)$. Integrating the independent
disturbances yields exactly the mediation functional above. This
derivation makes its stronger model explicit: splitting a node into
different edge values is not the same intervention as setting the
single node once. A manipulation of a composite product that changes
only a pricing channel also requires the claimed separation of those
mechanisms to have operational meaning.

Now introduce a post-action covariate $L$ that affects both $M$ and
$Y$. The natural response $Y(a,M(a'))$ involves $L(a)$ in the outcome
path and $L(a')$ inside the mediator path. The observed conditional
kernels do not generally determine the joint law of these two
counterfactual values. The same node would have to transmit different
versions of its history along its outgoing paths; ordinary
single-world g-computation does not supply that joint law.

A feasible alternative can be defined by a new randomized mediator
draw $G_{a'}$ from $p(m\mid A=a',X)$, conditionally independent of
the customer's other innovations given $X$. Under action
exchangeability, mediator--outcome exchangeability given $(A,L,X)$,
consistency and the requisite cross-regime support, direct mechanism
replacement identifies the project-defined target
\begin{equation}
 \E[Y(a,G_{a'})]=\sum_{x,l,m}
 \E[Y\mid a,l,m,x]p(l\mid a,x)p(m\mid a',x)p(x).
 \label{r3:random-mediator}
\end{equation}
The proof is the sequential g-formula: draw $X$, let $L$ develop
under $a$, independently draw the mediator from the stipulated
baseline-conditional distribution, then apply the invariant outcome
kernel. This intervention requires no joint $(L(a),L(a'))$ law.
It can be unrealistic if the drawn mediator is infeasible for the
resulting $L$, so support and operational feasibility remain essential.

Differences between $(a,a')$ pairs decompose effects of these new
interventions. They do not automatically decompose the original
total effect: even $Y(a,G_a)$ can differ from $Y(a)$ because the
fresh draw removes the natural association between $L$ and $M$.
For example, if $L$ is fair binary, $M=L$ and $Y=LM$, the natural
mean is $1/2$ while independently drawing $G$ from the same binary
margin gives $1/4$. This algebraic example explains the target
difference; it does not satisfy mediator support for estimating the
randomized intervention from that deterministic observational law.
Stochastic mediator policies conditional on the realized $L$ define
yet another target. A product must preserve the chosen mechanism
and conditioning information instead of returning all of these as
one ``indirect effect.''

\section{What the evidence allows us to conclude}
\label{r3:graph-contract}
The mathematical input is an intervention, target population, observed
information and model class. Its identification output is either a
functional with an assumption-dependent derivation, a set of possible
values, or a demonstrated failure of identification. A missing software
implementation is a fourth state: it says nothing about identification.

For the graph in Figure~\ref{r3:bank-graph}, a result cannot be labeled
``adjusted ATE'' after simply fitting $Y$ on $(A,X)$. A valid result might
instead report the randomized invitation effect, the consistency bounds,
or an IV analysis with extra assumptions. If multiple candidate graphs
remain, compute the effect for each supported graph or carry their union
of identified sets. Reporting the spread across two handpicked graphs is
only a sensitivity exercise for those graphs, not a confidence interval
covering all plausible causal structures.

The initial product can implement explicit DAG back-door checks and
registered functional templates with human-supplied assumptions. It should
record variable meanings, graph version, intervention transformation,
required kernels, support restrictions and target population weights.
Discovery and automated data preparation can propose inputs to this
process; they cannot certify the absence of unmeasured causes. The next
chapter shows why an instrumental-variable route changes the returned
object even when it uses the same customers, actions and outcomes.
